\section{Renormalization flow and physical correlations}

\begin{theorem}[Canonical-form blocks are injective]\label{thm:rfp_cf_structural}
    \lean{MPSTensor.rfp_cf_structural}
    \leanok
    \uses{def:is_canonical_form, def:injective}
    For a family of tensors in canonical form, every block is injective.
    This is part of the canonical-form hypotheses in
    \cite[Section~2.3]{Cirac2016MPDO_arXiv}.
\end{theorem}
\begin{proof}\leanok
    The canonical-form hypotheses include injectivity of each block.
\end{proof}


\begin{theorem}[RG flow convergence for canonical-form tensors]%
    \label{thm:rg_flow_converges_of_cf}
    \lean{MPSTensor.rg_flow_converges_of_cf}
    \leanok
    \uses{def:is_canonical_form, def:transfer_map, def:fixed_point_proj,
        thm:exponential_convergence_primitive}
    For a tensor in canonical form with blocks $(A_k)$ and weights $(\mu_k)$,
    the iterated blocking $\E_{A_k}^{2^n}$ converges pointwise to an
    idempotent linear map $\E_\infty$ as $n\to\infty$.
    That is, there exists $\E_\infty$ with $\E_\infty^2 = \E_\infty$ such that
    $\E_{A_k}^{2^n}(\rho) \to \E_\infty(\rho)$ for all density matrices $\rho$.
    This is \cite[Appendix~B]{Cirac2016MPDO_arXiv}.
\end{theorem}
\begin{proof}\leanok
    Each block is injective and left-canonical, so its transfer map is a
    primitive channel with a unique positive-definite fixed point $\rho_0$.
    The exponential convergence bound
    $\lVert \E^n(X) - P(X)\rVert \le C(1-\delta)^n \lVert X\rVert$
    (Theorem~\ref{thm:exponential_convergence_primitive}) then gives
    pointwise convergence $\E^n(X)\to P(X)$, and composing with the
    subsequence $2^n\to\infty$ yields the result.
    The witness $\E_\infty = P$ is the rank-one fixed-point projection,
    which is idempotent.
\end{proof}

\begin{remark}\label{rem:zero_correlation_length}
    \uses{thm:spectral_radius_le_one, def:correlation_length,
        def:mps_transfer_idempotence}
    When the transfer map is idempotent ($\E^2=\E$), connected correlations become
    separation-independent for all separations $n \ge 1$.
    Informally, this corresponds to zero correlation length ($\xi=0$).
    The full spectral equivalence (idempotence iff vanishing subleading spectrum)
    is deferred here.
\end{remark}

\begin{definition}[Virtual-insertion distance independence]\label{def:is_cid}
    \lean{MPSTensor.IsCID}
    \leanok
    \uses{def:connected_correlator}
    For a positive definite right fixed point $\rho^R > 0$, the virtual
    insertion expression is independent of distance when for all virtual bond
    matrices $X$ and $Y$ and all separations $n,m \ge 1$,
    \[
        C(X,Y;n) = C(X,Y;m).
    \]
\end{definition}

\begin{definition}[Physical observable transfer]
    \label{def:physical_observable_transfer}
    \lean{MPSTensor.physicalObservableTransfer}
    \leanok
    \uses{def:eval_word}
    For an observable $O$ on a block of $L$ physical spins, define the linear
    map on the virtual matrix space by
    \[
        \E_O(X)=\sum_{\boldsymbol i,\boldsymbol j}
        \langle\boldsymbol j|O|\boldsymbol i\rangle
        A^{\boldsymbol i}X(A^{\boldsymbol j})^\dagger .
    \]
    This is the observable transfer map of
    \cite[lines~490--496]{Cirac2016MPDO_arXiv}.
\end{definition}

\begin{definition}[Two-observable periodic expectation]
    \label{def:physical_two_point_expectation}
    \lean{MPSTensor.physicalTwoPointExpectation}
    \leanok
    \uses{def:physical_observable_transfer, def:transfer_map}
    Let $O_1,O_2$ act on blocks of $L_1,L_2$ physical spins, separated around
    the periodic chain by complementary gaps of lengths $n_1,n_2$.  Their
    two-observable expectation is
    \[
        \Tr_{\MN{D}}\!\left(
        \E_{O_2}\circ\E_A^{n_2}\circ\E_{O_1}\circ\E_A^{n_1}
        \right).
    \]
    The trace is the operator trace on the virtual matrix space, as in
    \cite[lines~490--496]{Cirac2016MPDO_arXiv}.
\end{definition}

\begin{definition}[Physical correlations independent of distance]%
    \label{def:source_physical_cid}
    \lean{MPSTensor.IsPhysicalCID}
    \leanok
    \uses{def:physical_two_point_expectation}
    Let $O_1,O_2$ act on two disjoint contiguous regions of positive lengths
    $L_1,L_2\geq1$ in a periodic chain, with complementary gap lengths
    $n_1,n_2\geq0$. Physical correlations are independent of distance when, for
    every $m_1,m_2\geq0$ satisfying
    $n_1+n_2=m_1+m_2$, one has
    \[
      \Tr_{\MN{D}}\!\left(
        \E_{O_2}\circ\E_A^{n_2}\circ\E_{O_1}\circ\E_A^{n_1}\right)
      =
      \Tr_{\MN{D}}\!\left(
        \E_{O_2}\circ\E_A^{m_2}\circ\E_{O_1}\circ\E_A^{m_1}\right).
    \]
    Thus either region may be translated without crossing the other. Zero
    gaps, corresponding to adjacent regions, are included. This is
    \cite[Definition~3.3 and lines~490--496]{Cirac2016MPDO_arXiv}.
\end{definition}

\begin{definition}[Positive-gap physical distance independence]%
    \label{def:is_positive_gap_physical_cid}
    \lean{MPSTensor.IsPositiveGapPhysicalCID}
    \leanok
    \uses{def:physical_observable_transfer,
        def:physical_two_point_expectation}
    For an observable $O$ on a block of $L$ physical spins, write
    \[
        \E_O(X)=\sum_{\boldsymbol i,\boldsymbol j}
        \langle\boldsymbol j|O|\boldsymbol i\rangle
        A^{\boldsymbol i}X(A^{\boldsymbol j})^\dagger .
    \]
    Place two observables $O_1,O_2$ on nonempty finite blocks of a periodic
    chain, with positive complementary gap lengths $n_1,n_2$. The tensor has
    \emph{positive-gap physical correlations independent of distance} when the expectation
    is unchanged upon replacing these gaps by positive $m_1,m_2$ satisfying
    $n_1+n_2=m_1+m_2$:
    \[
      \Tr_{\MN{D}}\!\left(
        \E_{O_2}\circ\E_A^{n_2}\circ\E_{O_1}\circ\E_A^{n_1}\right)
      =
      \Tr_{\MN{D}}\!\left(
        \E_{O_2}\circ\E_A^{m_2}\circ\E_{O_1}\circ\E_A^{m_1}\right).
    \]
    This is the transfer-matrix form of
    \cite[Definition~3.3 and the correlation formula preceding
    Theorem~3.8]{Cirac2016MPDO_arXiv}, restricted to positive complementary
    gaps.  It excludes adjacent regions.
\end{definition}

\begin{lemma}[Physical CID implies positive-gap physical CID]%
    \label{lem:is_positive_gap_physical_cid_of_physical_cid}
    \lean{MPSTensor.isPositiveGapPhysicalCID_of_isPhysicalCID}
    \leanok
    \uses{def:source_physical_cid, def:is_positive_gap_physical_cid}
    Correlations independent of distance for all disjoint physical regions are
    independent of distance when both complementary gaps are positive.
\end{lemma}
\begin{proof}\leanok
    Restrict the four gap lengths to positive integers.
\end{proof}

\begin{theorem}[Idempotent transfer implies positive-gap physical CID]%
    \label{thm:is_positive_gap_physical_cid_of_is_rfp}
    \lean{MPSTensor.isPositiveGapPhysicalCID_of_isTransferIdempotent}
    \leanok
    \uses{def:mps_transfer_idempotence, def:is_positive_gap_physical_cid}
    If $\E_A^2=\E_A$, then physical correlations are independent of distance
    whenever both complementary gaps are positive.
\end{theorem}
\begin{proof}\leanok
    Idempotence gives $\E_A^n=\E_A$ for every $n\geq1$.  Substituting
    $\E_A^{n_1}=\E_A^{n_2}=\E_A^{m_1}=\E_A^{m_2}=\E_A$ in the two-observable
    transfer formula gives
    \[
      \Tr_{\MN{D}}\!\left(
        \E_{O_2}\circ\E_A^{n_2}\circ\E_{O_1}\circ\E_A^{n_1}\right)
      =
      \Tr_{\MN{D}}\!\left(
        \E_{O_2}\circ\E_A\circ\E_{O_1}\circ\E_A\right)
      =
      \Tr_{\MN{D}}\!\left(
        \E_{O_2}\circ\E_A^{m_2}\circ\E_{O_1}\circ\E_A^{m_1}\right).
    \]
\end{proof}

\begin{definition}[Zero correlation length]\label{def:is_zcl}
    \lean{MPSTensor.IsZCL}
    \leanok
    \uses{def:is_locally_orthogonal, def:is_bnt_zcl, def:is_cid}
    In the single-block convention, an MPS tensor has \emph{zero correlation
        length} when
    \[
        \E_A^2=\E_A
        \qquad\text{and}\qquad
        \operatorname{CID}(A)
    \]
    both hold. The source definition is
    Definition~\ref{def:source_bnt_zcl}; the present condition uses virtual
    insertions instead.
\end{definition}

\begin{theorem}[Virtual-insertion distance independence implies RFP]%
    \label{thm:isCID_implies_isTransferIdempotent}
    \lean{MPSTensor.isCID_implies_isTransferIdempotent}
    \leanok
    \uses{def:is_cid, def:mps_transfer_idempotence, def:transfer_map}
    If an MPS tensor admits a positive definite right fixed point and its
    virtual-insertion correlator is independent of distance, then its transfer
    map is idempotent.
    % The reverse direction of Theorem 3.8 in arXiv:1606.00608 additionally
    % fails for raw repeated-copy weights. A corrected restricted theorem would
    % additionally require realizing virtual insertions by physical observables.
\end{theorem}
\begin{proof}\leanok
    Let $\rho>0$ be the fixed point.  For every matrix $Z$, invertibility gives
    $Z=(Z\rho^{-1})\rho$.  Apply distance independence with
    $X=Z\rho^{-1}$, the second insertion $N$, and separations $2$ and $1$.
    The identical disconnected terms cancel.  Thus, for every $N$,
    \[
        \tr\!\left(N\E_A^2(Z)\right)
        =\tr\!\left(N\E_A(Z)\right).
    \]
    Nondegeneracy of the trace pairing yields
    $\E_A^2(Z)=\E_A(Z)$ for every $Z$, hence $\E_A^2=\E_A$.
\end{proof}

\begin{theorem}[Zero correlation length iff idempotent transfer]%
    \label{thm:zcl_iff_idempotent_transfer}
    \lean{MPSTensor.zcl_iff_idempotent_transfer}
    \leanok
    \uses{def:mps_transfer_idempotence, def:is_cid, def:is_locally_orthogonal}
    Under the single-block convention above,
    \[
        \bigl(\E_A^2=\E_A \ \text{and}\ \operatorname{CID}(A)\bigr)
        \quad\Longleftrightarrow\quad
        \E_A^2=\E_A .
    \]
    This is the idempotence/CID part of
    \cite[Theorem~3.8]{Cirac2016MPDO_arXiv}. The full source theorem for
    canonical-form tensors instead uses the unrestricted physical condition in
    Definition~\ref{def:source_bnt_zcl}; under its stated raw-weight
    normalization, that biconditional is false by
    Theorem~\ref{thm:halved_weight_zcl_not_rfp}.
\end{theorem}
\begin{proof}\leanok
    The forward direction extracts the idempotence hypothesis.  Conversely,
    idempotence gives $\E_A^n=\E_A$ for every $n\geq 1$.  Hence, for every
    positive-definite right fixed point $\rho^R$, every pair of virtual
    matrices $X,Y$, and every $n\geq 1$,
    \[
        \begin{aligned}
        C(X,Y;n)
        &=\tr\!\left(Y\E_A^n(X\rho^R)\right)
          -\tr(X\rho^R)\tr(Y\rho^R)\\
        &=\tr\!\left(Y\E_A(X\rho^R)\right)
          -\tr(X\rho^R)\tr(Y\rho^R)\\
        &=C(X,Y;1).
        \end{aligned}
    \]
    Thus the connected correlator is independent of every positive
    separation.
\end{proof}

\begin{theorem}[Single-block RFP iff ZCL inside a canonical form]%
    \label{thm:rfp_iff_zcl}
    \lean{MPSTensor.rfp_iff_zcl}
    \leanok
    \uses{def:is_canonical_form, def:mps_transfer_idempotence, def:is_zcl}
    Let $A^i=\bigoplus_k \mu_k A_k^i$ be a tensor in canonical form. For each
    $k$, the tensor $A_k$ is a renormalization fixed point if and only if
    $A_k$ has zero correlation length.

    This is the single-block consequence of
    \cite[Theorem~3.10]{Cirac2016MPDO_arXiv}. The ZCL condition in the source
    theorem is the BNT-family condition: CID together with the local
    orthogonality equations $\mathbb E_{j,j'}=0$ for all $j\ne j'$,
    as in Definition~\ref{def:source_bnt_zcl}.
\end{theorem}
\begin{proof}\leanok
    \uses{thm:zcl_iff_idempotent_transfer}
    Apply Theorem~\ref{thm:zcl_iff_idempotent_transfer} to the chosen
    canonical-form block.
\end{proof}
